% Generated draft: requires mathematical and source review.
\documentclass[11pt]{article}
\usepackage[margin=25mm]{geometry}
\usepackage{fontspec}
\setmainfont{Latin Modern Roman}
\usepackage{amsmath,amssymb,mathtools,unicode-math}
\setmathfont{Latin Modern Math}
\usepackage{enumitem,needspace,xurl,hyperref}
\setlength{\parindent}{0pt}
\setlength{\parskip}{5pt}
\setlength{\emergencystretch}{3em}
\newcommand{\N}{\mathbb{N}}
\newcommand{\Z}{\mathbb{Z}}
\newcommand{\Q}{\mathbb{Q}}
\newcommand{\R}{\mathbb{R}}
\newcommand{\C}{\mathbb{C}}
\newcommand{\F}{\mathbb{F}}
\newcommand{\A}{\mathbb{A}}
\newcommand{\PP}{\mathbb{P}}
\newcommand{\E}{\mathbb{E}}
\newcommand{\eps}{\varepsilon}
\newcommand{\dd}{\,\mathrm{d}}
\begin{document}
\section*{UnsolvedMath ID: 701763}
\textbf{Source problem code:} catalogue entry
\subsection{Definitions and mathematical statement}
Let \(p\) be a prime. A group \(G\) is a \(p\)-group if every element has finite order equal to a power of \(p\). Its commutator subgroup is
\[
[G,G]=\langle x^{-1}y^{-1}xy:x,y\in G\rangle,
\]
and \(G\) is perfect if \(G=[G,G]\).

For any group \(H\), define its transfinite upper central series by
\[
Z_0(H)=\{1\},\qquad
Z_{\alpha+1}(H)/Z_\alpha(H)=Z\bigl(H/Z_\alpha(H)\bigr),\qquad
Z_\lambda(H)=\bigcup_{\alpha<\lambda}Z_\alpha(H)
\]
for every ordinal \(\alpha\) and every nonzero limit ordinal \(\lambda\), where \(Z(K)\) denotes the center of a group \(K\). The group \(H\) is nilpotent if \(Z_n(H)=H\) for some nonnegative integer \(n\), and hypercentral if \(Z_\alpha(H)=H\) for some ordinal \(\alpha\). A group is locally nilpotent if each of its finitely generated subgroups is nilpotent.

The supplied question asks whether there exist a prime \(p\) and a group \(G\) such that
\[
\begin{gathered}
G\text{ is a perfect, locally nilpotent }p\text{-group},\\
\forall H<G\;\exists\text{ an ordinal }\alpha_H\quad Z_{\alpha_H}(H)=H.
\end{gathered}
\]
Here \(H<G\) means any proper subgroup, without a normality or finite-generation restriction, and \(\alpha_H\) may depend on \(H\).

The source does not explicitly require \(G\ne\{1\}\). Under the definitions above, the trivial group is permitted by the stated conditions. Whether a nontrivial group is intended is therefore an unresolved qualification; the additional hypothesis \(G\ne\{1\}\) has not been silently inserted.
\subsection{Short English statement}
Does there exist, for some prime \(p\), a perfect locally nilpotent \(p\)-group all of whose proper subgroups are hypercentral? The source leaves nontriviality unstated, so the intended requirement concerning \(G\ne\{1\}\) needs clarification.
\subsection{Sources}
\begin{itemize}
\item Supplied local UnsolvedMath registry record, permanent ID \(701763\), exact source code \texttt{catalogue entry}, titled \emph{Kourovka Notebook Problem 20.105}. The registry labels the problem open; this is supplied source metadata, not an independently verified assessment.
\item \emph{The Kourovka Notebook (Unsolved Problems in Group Theory)}, 20th issue (2022), Problem 20.105; the supplied background also identifies the 21st issue edition (2026), arXiv:1401.0300v48. The collection description names E. I. Khukhro and V. D. Mazurov as editors. These bibliographic details are supplied and have not been independently checked.
\item Proposed by N. Trabelsi; proposed year 2022, according to the supplied record.
\end{itemize}
\end{document}
